% kp.tex — dikonversi dari makalah-hadis/draft/KATA_PENGANTAR.md (tanpa heading)

Puji syukur kehadirat Allah SWT atas rahmat, hidayah, dan inayah-Nya sehingga penulis dapat menyelesaikan penyusunan makalah yang berjudul \textbf{\textit{Dari Sahifah ke Basis Data: Periodisasi Kodifikasi Hadis, Penanggulangan Fabrikasi, dan Tantangan Otoritas di Era Digital}} ini dengan lancar. Shalawat serta salam semoga senantiasa tercurahkan kepada junjungan kita Nabi Agung Muhammad SAW, keluarga, sahabat, hingga umatnya yang konsisten menjaga kemurnian risalah kenabian hingga akhir zaman. Pemilihan topik ini didasari oleh dinamika krusial dalam studi hadis kontemporer, di mana pergeseran media transmisi---dari catatan personal (\textit{sahifah}) klasik menuju basis data dan algoritma digital---membawa peluang sekaligus ancaman epistemologis. Di satu sisi, teknologi mempercepat daya jangkau dan aksesibilitas penelusuran riwayat; di sisi lain, fenomena viralitas tanpa verifikasi, merosotnya literasi kritik, dan maraknya disinformasi agama di ruang siber menuntut rekonstruksi historiografis yang padu mengenai periodisasi perkembangan hadis, arsitektur penyelamatan dari pemalsuan (\textit{maudu'}), serta mediatisasi otoritas keagamaan di era modern.

Makalah ini disusun guna memenuhi tugas akademik mata kuliah Ulumul Hadis pada Program Studi Ilmu Falak, Fakultas Syariah dan Hukum, Universitas Islam Negeri Walisongo Semarang tahun akademik 2026/2027, sekaligus diharapkan memberikan kontribusi pemikiran dalam menjembatani khazanah filologi klasik dan tantangan verifikasi digital. Dalam proses penyusunannya, penulis mendapatkan bimbingan, arahan, dan dorongan dari berbagai pihak. Oleh karena itu, penulis menyampaikan rasa terima kasih dan penghargaan yang setinggi-tingginya kepada Bapak Prof. DR. H. Nama Dosen, M.S.I., selaku dosen pengampu mata kuliah Ulumul Hadis, yang telah memberikan petunjuk keilmuan serta wawasan metodologis yang sangat berharga. Ucapan terima kasih juga penulis sampaikan kepada rekan-rekan mahasiswa Program Studi Ilmu Falak serta seluruh pihak yang telah memberikan bantuan, diskusi kritis, dan dukungan moral selama proses penyelesaian karya tulis ini.

Penulis menyadari sepenuhnya bahwa makalah ini masih memiliki keterbatasan dan kekurangan, baik dari segi kedalaman analisis maupun jangkauan pembahasan. Oleh sebab itu, penulis secara terbuka mengharapkan kritik konstruktif serta saran yang membangun dari para pembaca demi penyempurnaan karya-karya akademis di masa mendatang. Akhir kata, semoga makalah ini dapat memberikan kemanfaatan akademis yang nyata, memperkaya literasi keilmuan, serta menjadi sumbangsih pemikiran yang berguna bagi perkembangan studi hadis dan khazanah keislaman pada umumnya.

\vspace{1.5cm}

\noindent
\begin{minipage}[t]{0.45\textwidth}
    \centering
    Mengetahui,\\
    Dosen Pengampu,

    \vspace{2.5cm}

    \textbf{Prof. DR. H. Nama Dosen,}\\
    \textbf{M.S.I.}
\end{minipage}%
\hfill
\begin{minipage}[t]{0.45\textwidth}
    \centering
    Semarang, 25 September 2026\\
    Penulis,

    \vspace{2.5cm}

    \textbf{Nama Penulis 1}\\
    \textbf{Nama Penulis 3}
\end{minipage}

